% !TeX program = lualatex
% =====================================================================
%  Operador de Sturm–Liouville y ecuación de Hermite
%  Apuntes formalizados (estilo oscuro, a juego con «Mecánica I»)
% =====================================================================
\documentclass[11pt,a4paper,twoside]{article}

% ---------- Datos editables ------------------------------------------
\newcommand{\autor}{Daniel Martín Guerrero}
\newcommand{\asignatura}{Métodos matemáticos}
\newcommand{\titulodoc}{Operador de Sturm--Liouville y ecuación de Hermite}
\newcommand{\titulocorto}{Sturm--Liouville y Hermite}

% ---------- Paquetes --------------------------------------------------
\usepackage[spanish,es-noshorthands,es-nolists,es-tabla,es-lcroman]{babel}
\usepackage{fontspec}
\usepackage{amsmath,amssymb,mathtools,bm,esint}
\usepackage[a4paper,top=2.9cm,bottom=2.7cm,left=2.6cm,right=2.6cm,
            headsep=1.0cm,headheight=14pt]{geometry}
\usepackage{xcolor}
\usepackage{fancyhdr}
\usepackage{titlesec}
\usepackage{enumitem}
\usepackage[most]{tcolorbox}
\usepackage{multicol}
\usepackage{booktabs,array}
\usepackage{tikz}
\usetikzlibrary{arrows.meta,calc,positioning,decorations.pathreplacing}
\usepackage{pgfplots}
\pgfplotsset{compat=1.18}
\usepgfplotslibrary{fillbetween}
\usepackage{caption}
\usepackage{xparse}
\usepackage{hyperref}

% ---------- Paleta ----------------------------------------------------
\definecolor{bg}{HTML}{000000}
\definecolor{fg}{HTML}{ECECEC}
\definecolor{boxbg}{HTML}{0E1014}
\definecolor{acc}{HTML}{5AA9E6}
\definecolor{accdim}{HTML}{3D7198}
\definecolor{lav}{HTML}{8FA3D6}
\definecolor{gold}{HTML}{D6B25E}
\definecolor{muted}{HTML}{9CA0A8}
\definecolor{rulec}{HTML}{2E3036}
\definecolor{sky}{HTML}{B7CCE8}
% gradiente de niveles
\definecolor{c0}{HTML}{4F7FE0}
\definecolor{c1}{HTML}{3FC1C9}
\definecolor{c2}{HTML}{8FD14F}
\definecolor{c3}{HTML}{E6C24A}
\definecolor{c4}{HTML}{F08A3C}
\definecolor{c5}{HTML}{E0452F}

\pagecolor{bg}
\AtBeginDocument{\color{fg}}

\hypersetup{
  colorlinks=true, linkcolor=acc, urlcolor=acc, citecolor=acc,
  pdftitle={\titulodoc}, pdfauthor={\autor},
  pdfsubject={Sturm--Liouville, Hermite}
}

% ---------- Cabeceras -------------------------------------------------
\pagestyle{fancy}
\fancyhf{}
\fancyhead[LO]{\normalfont\small\color{muted}\textsc{\titulocorto}}
\fancyhead[RO]{\normalfont\small\color{muted}\thepage}
\fancyhead[LE]{\normalfont\small\color{muted}\thepage}
\fancyhead[RE]{\normalfont\small\color{muted}\textsc{Apuntes formalizados}}
\makeatletter
\renewcommand{\headrule}{{\color{rulec}\hrule\@height0.4pt\@width\headwidth\vskip-0.4pt}}
\makeatother
\renewcommand{\footrulewidth}{0pt}

% ---------- Secciones -------------------------------------------------
\titleformat{\section}{\Large\bfseries}{\color{acc}\thesection}{0.9em}{}
  [\vspace{-0.4ex}{\color{rulec}\titlerule[0.6pt]}]
\titleformat{\subsection}{\large\bfseries}{\color{acc}\thesubsection}{0.9em}{}
\titlespacing*{\section}{0pt}{3.2ex plus 1ex}{2.2ex}
\titlespacing*{\subsection}{0pt}{2.6ex plus 1ex}{1.4ex}
\numberwithin{equation}{section}

\setlist[enumerate]{label=\textcolor{acc}{\arabic*.},leftmargin=1.8em,itemsep=0.35ex}
\setlist[itemize]{label=\textcolor{acc}{\scriptsize$\blacktriangleright$},leftmargin=1.6em,itemsep=0.3ex}

\DeclareCaptionFont{mutedc}{\color{muted}}
\DeclareCaptionFont{accc}{\color{acc}}
\captionsetup{font={small,mutedc},labelfont={sc,accc},labelsep=period,
              justification=justified,margin=1.2cm}

% ---------- Entornos formalizados --------------------------------------
\newcounter{thm}[section]
\renewcommand{\thethm}{\thesection.\arabic{thm}}

\tcbset{
  thmbox/.style={enhanced, breakable,
    colback=boxbg, colframe=#1, coltext=fg, boxrule=0.5pt, arc=1.2pt,
    left=4mm, right=4mm, top=4.2mm, bottom=3mm,
    before skip=3.2ex, after skip=2.6ex,
    attach boxed title to top left={xshift=4mm,yshift*=-\tcboxedtitleheight/2},
    boxed title style={colback=boxbg, colframe=#1, boxrule=0.5pt, arc=1.2pt,
                       left=2mm,right=2mm,top=0.6mm,bottom=0.6mm},
    coltitle=#1, fonttitle=\normalfont},
}

\NewDocumentEnvironment{definicion}{o}
 {\refstepcounter{thm}\begin{tcolorbox}[thmbox=acc,
   title={\textsc{Definición \thethm}\IfValueT{#1}{\ \textit{(#1)}}}]}
 {\end{tcolorbox}}
\NewDocumentEnvironment{proposicion}{o}
 {\refstepcounter{thm}\begin{tcolorbox}[thmbox=lav,
   title={\textsc{Proposición \thethm}\IfValueT{#1}{\ \textit{(#1)}}}]\itshape}
 {\end{tcolorbox}}
\NewDocumentEnvironment{corolario}{o}
 {\refstepcounter{thm}\begin{tcolorbox}[thmbox=lav,
   title={\textsc{Corolario \thethm}\IfValueT{#1}{\ \textit{(#1)}}}]\itshape}
 {\end{tcolorbox}}
\NewDocumentEnvironment{teorema}{o}
 {\refstepcounter{thm}\begin{tcolorbox}[thmbox=gold,
   title={\textsc{Teorema \thethm}\IfValueT{#1}{\ \textit{(#1)}}}]\itshape}
 {\end{tcolorbox}}

\NewDocumentEnvironment{observacion}{o}
 {\par\smallskip\refstepcounter{thm}\noindent
  {\color{lav}\bfseries\itshape Observación \thethm}%
  \IfValueT{#1}{\ {\color{muted}(#1)}}.\ }
 {\par\smallskip}
\NewDocumentEnvironment{ejemplo}{o}
 {\par\smallskip\refstepcounter{thm}\noindent
  {\color{acc}\bfseries Ejemplo \thethm}%
  \IfValueT{#1}{\ {\color{muted}(#1)}}.\ }
 {\par\smallskip}
\NewDocumentEnvironment{demostracion}{o}
 {\par\smallskip\noindent{\bfseries\itshape Demostración}%
  \IfValueT{#1}{\ \textit{#1}}.\ }
 {\hfill{\color{muted}$\square$}\par\smallskip}

\newtcolorbox{metodo}[1]{enhanced, breakable, colback=boxbg, coltext=fg,
  frame hidden, sharp corners, borderline west={2.2pt}{0pt}{acc},
  left=5mm, right=4mm, top=2.5mm, bottom=2.5mm, before skip=2.6ex, after skip=2.6ex,
  before upper={{\color{acc}\textsc{Método: #1.}}\par\smallskip}}

% resaltado de resultados
\newcommand{\res}[1]{\tcbhighmath[colback=boxbg,colframe=acc,colupper=fg,boxrule=0.5pt,arc=1pt,
  left=1.5mm,right=1.5mm,top=1mm,bottom=1mm]{#1}}

% ---------- Notación --------------------------------------------------
\newcommand{\dd}{\mathrm{d}}
\newcommand{\dv}[2]{\frac{\dd #1}{\dd #2}}
\newcommand{\dvn}[3]{\frac{\dd^{#3} #1}{\dd #2^{#3}}}
\newcommand{\pd}{\partial}
\newcommand{\Lop}{\mathcal{L}}
\newcommand{\R}{\mathbb{R}}
\newcommand{\C}{\mathbb{C}}
\newcommand{\N}{\mathbb{N}}
\newcommand{\Z}{\mathbb{Z}}
\newcommand{\Ap}{A_{+}}
\newcommand{\Am}{A_{-}}
\newcommand{\ip}[2]{\left(#1,#2\right)}
\newcommand{\norm}[1]{\left\lVert #1\right\rVert}
\newcommand{\Ltwo}{L^{2}_{\rho}}
\newcommand{\defeq}{\stackrel{\mathrm{def}}{=}}
\newcommand{\e}{\mathrm{e}}
\newcommand{\iu}{\mathrm{i}}

% ---------- Figura: oscilador armónico --------------------------------
% \oscilador{anchura}{ejes? (true/false)}
\newcommand{\curvaHermite}[3]{% n, color, expresión de H_n
  \addplot[name path=lvl#1, draw=none, domain=-4.6:4.6] {#1+0.5};
  \addplot[name path=psi#1, #2, line width=0.9pt, samples=220, domain=-4.6:4.6]
     {#1+0.5+0.95*\csname N#1\endcsname*(#3)*exp(-x^2/2)};
  \addplot[#2, fill opacity=0.16] fill between[of=psi#1 and lvl#1];
  \draw[#2, opacity=0.35, line width=0.3pt]
     (axis cs:{-sqrt(2*#1+1)},#1+0.5) -- (axis cs:{sqrt(2*#1+1)},#1+0.5);
}
\expandafter\def\csname N0\endcsname{0.75113}
\expandafter\def\csname N1\endcsname{0.53113}
\expandafter\def\csname N2\endcsname{0.26556}
\expandafter\def\csname N3\endcsname{0.10842}
\expandafter\def\csname N4\endcsname{0.03833}
\expandafter\def\csname N5\endcsname{0.01212}

\newcommand{\oscilador}[2]{%
\begin{tikzpicture}
\begin{axis}[
  width=#1, height=0.62*#1,
  xmin=-4.6, xmax=4.6, ymin=-0.15, ymax=6.6,
  axis lines=#2, axis line style={muted, line width=0.4pt},
  tick style={muted}, ticklabel style={muted, font=\scriptsize},
  xtick={-4,-2,0,2,4}, ytick={0.5,1.5,2.5,3.5,4.5,5.5},
  yticklabels={$\tfrac12$,$\tfrac32$,$\tfrac52$,$\tfrac72$,$\tfrac92$,$\tfrac{11}{2}$},
  xlabel={\color{muted}$\xi$}, ylabel={\color{muted}$E/\hbar\omega$},
  x label style={at={(axis description cs:1.0,0.0)},anchor=north west},
  y label style={at={(axis description cs:0.5,1.0)},anchor=south,rotate=-90},
  clip=true]
  \addplot[sky, line width=0.7pt, samples=120, domain=-3.6:3.6] {x^2/2};
  \curvaHermite{0}{c0}{1}
  \curvaHermite{1}{c1}{2*x}
  \curvaHermite{2}{c2}{4*x^2-2}
  \curvaHermite{3}{c3}{8*x^3-12*x}
  \curvaHermite{4}{c4}{16*x^4-48*x^2+12}
  \curvaHermite{5}{c5}{32*x^5-160*x^3+120*x}
\end{axis}
\end{tikzpicture}}

\setlength{\parindent}{0pt}
\setlength{\parskip}{0.55ex}

% =====================================================================
\begin{document}

% ------------------------------ Portada ------------------------------
\begin{titlepage}
\newgeometry{top=2.2cm,bottom=2.2cm,left=2.6cm,right=2.6cm}
\centering
{\Large\scshape Universidad de Salamanca\par}
\vspace{0.9cm}
{\color{fg}\rule{\linewidth}{0.4pt}}\par\vspace{0.45cm}
{\small\color{muted}\scshape Facultad de Ciencias \quad$\cdot$\quad Grado en Física\par}
\vspace{1.4cm}
{\color{acc}\scshape \asignatura\par}
\vspace{0.15cm}
{\small\color{muted}\scshape Apuntes formalizados\par}
\vspace{1.0cm}
{\fontsize{24}{30}\selectfont\bfseries Operador de Sturm--Liouville\\[2pt] y ecuación de Hermite\par}
\vspace{0.7cm}
{\color{fg}\rule{0.2\linewidth}{0.4pt}}\par\vspace{0.4cm}
{\small\itshape Separación de variables \;$\cdot$\; Espacio de Hilbert $L^2_\rho$ \;$\cdot$\;
 Series, Rodrigues\\ y operadores escalera\par}
\vspace{0.8cm}
\oscilador{0.92\linewidth}{none}
\vspace{0.5cm}

{\small
$\displaystyle H_n(x)=(-1)^n\,\e^{x^2}\dvn{}{x}{n}\e^{-x^2},\qquad
\int_{-\infty}^{\infty}\!H_nH_m\,\e^{-x^2}\dd x=2^n n!\sqrt{\pi}\,\delta_{nm},\qquad
[\Am,\Ap]=2$\par}
\vfill
{\color{fg}\rule{\linewidth}{0.4pt}}\par\vspace{0.3cm}
\begin{minipage}[t]{0.32\linewidth}\raggedright
{\scriptsize\color{acc}\scshape Autor}\\[2pt]{\small\autor}
\end{minipage}\hfill
\begin{minipage}[t]{0.32\linewidth}\centering
{\scriptsize\color{acc}\scshape Contenido}\\[2pt]{\small Sturm--Liouville $\cdot$ Hermite}\\
{\scriptsize\color{muted}Desarrollo simbólico completo}
\end{minipage}\hfill
\begin{minipage}[t]{0.32\linewidth}\raggedleft
{\scriptsize\color{acc}\scshape Apuntes formalizados}\\[2pt]{\small Salamanca, curso 2026/2027}
\end{minipage}
\restoregeometry
\end{titlepage}

% ------------------------------ Índice -------------------------------
\setcounter{page}{2}
{\hypersetup{linkcolor=fg}\tableofcontents}
\vspace{1ex}
{\color{muted}\small Notación: $y'\equiv\dv{y}{x}$, $\;\ip{\cdot}{\cdot}$ producto escalar con peso,
$\;\Ltwo\equiv L^2_\rho$, $\;\N=\{0,1,2,\dots\}$, $\;\square$ fin de demostración.}

% =====================================================================
\section{Separación de variables y problema de autovalores}
% =====================================================================

\subsection{Reducción a un problema de autovalores}

\begin{definicion}[Problema de partida]
Sea $y(x,t)$, $x\in[a,b]$, $t\in\R$, solución de
\begin{equation}\label{eq:edp}
k(x)\,\pd_t^2 y=\pd_x\!\bigl(u(x)\,\pd_x y\bigr)-M(x)\,y,
\qquad k>0,\;u>0 .
\end{equation}
\end{definicion}

Ansatz de variables separadas, $y(x,t)=X(x)\,T(t)$:
\begin{align}
k\,X\,\ddot T &= T\,(uX')'-M\,X\,T
\quad\Big/\;\div\,(kXT)\notag\\
\frac{\ddot T(t)}{T(t)} &= \frac{(uX')'-MX}{k\,X}\Big|_{x}\equiv-\lambda,
\qquad \lambda=\text{const.}\label{eq:sep}
\end{align}
(miembro izquierdo sólo de $t$, derecho sólo de $x$ $\Rightarrow$ ambos constantes). Por tanto
\begin{equation}\label{eq:SLparcial}
\dvn{T}{t}{2}=-\lambda\,T(t),\qquad
\underbrace{\frac{1}{k}\Bigl[\dv{}{x}\Bigl(-u(x)\dv{X}{x}\Bigr)+M(x)\,X\Bigr]}_{\text{operador de Sturm--Liouville}}=\lambda\,X .
\end{equation}

Parte temporal:
\begin{equation}\label{eq:T}
\ddot T=-\lambda T\;\Longrightarrow\;T(t)=c_{+}\,\e^{\iu\omega t}+c_{-}\,\e^{-\iu\omega t},
\qquad \omega\defeq\sqrt{\lambda}.
\end{equation}
$\lambda$ queda fijado por las condiciones de contorno sobre $X$.

\begin{observacion}[Notación]
En el manuscrito, $T=\e^{\iu\lambda t}$; con \eqref{eq:T} la frecuencia es $\omega=\sqrt\lambda$
(ambas coinciden etiquetando el espectro por $\omega$).
\end{observacion}

\begin{ejemplo}[Cuerda homogénea con extremos fijos]
$k=u=1$, $M=0$, $[a,b]=[0,\pi]$, $X(0)=X(\pi)=0$:
\begin{enumerate}
\item $-X''=\lambda X$. Multiplicando por $X$ e integrando por partes:
$\lambda\int_0^\pi X^2\dd x=\int_0^\pi X'^2\dd x\ge0$; $\;\lambda=0\Rightarrow X'=0\Rightarrow X=0$. Luego $\lambda>0$.
\item $X=A\sin\sqrt\lambda\,x+B\cos\sqrt\lambda\,x$; $\;X(0)=0\Rightarrow B=0$;
$\;X(\pi)=0\Rightarrow\sin\sqrt\lambda\,\pi=0$.
\item Espectro discreto (cuantizado):
\begin{equation}
\res{\sqrt{\lambda_n}=\omega_n=n,\qquad \lambda_n=n^2,\qquad X_n(x)=\sin nx,\qquad n\in\N\setminus\{0\}.}
\end{equation}
\item Solución general (serie de Fourier):
\begin{equation}\label{eq:fourier}
y(x,t)=\sum_{n=1}^{\infty}\bigl(a_n\cos nt+b_n\sin nt\bigr)\sin nx .
\end{equation}
\item Con $\int_0^\pi\sin nx\,\sin mx\,\dd x=\tfrac{\pi}{2}\delta_{nm}$:
\[
a_n=\frac{2}{\pi}\int_0^\pi y(x,0)\sin nx\,\dd x,\qquad
b_n=\frac{2}{n\pi}\int_0^\pi \pd_t y(x,0)\sin nx\,\dd x .
\]
\end{enumerate}
$\{X_n\}$: autovectores del operador diferencial $\equiv$ base del espacio vectorial
$\Rightarrow$ análisis armónico de señales periódicas.
\end{ejemplo}

% =====================================================================
\section{Operador de Sturm--Liouville}
% =====================================================================

\subsection{Definición y producto escalar}

\begin{definicion}[Operador de Sturm--Liouville]
Sean $\rho,p\in\mathcal C^1(a,b)$, $q\in\mathcal C^0(a,b)$ reales, $\rho>0$ (\emph{función peso}), $p>0$:
\begin{equation}\label{eq:SL}
\Lop y\defeq\frac{1}{\rho(x)}\Bigl[-\dv{}{x}\Bigl(p(x)\dv{y}{x}\Bigr)+q(x)\,y(x)\Bigr],
\qquad y\in D_{\Lop}.
\end{equation}
Dominio: $D_{\Lop}=\{y\in\mathcal C^2(a,b):\ \mathcal B_a y=\mathcal B_b y=0\}$ (condiciones de contorno).
Problema de autovalores: $\Lop y=\lambda y$, $y\in D_{\Lop}\setminus\{0\}$ $\;\leadsto\;$ terna $(D_{\Lop},\lambda,y_\lambda)$.
\end{definicion}

Comparando \eqref{eq:SLparcial} con \eqref{eq:SL}: $\rho=k$, $p=u$, $q=M$.

\begin{definicion}[Espacio de Hilbert $L^2_\rho$]
\begin{gather}
\ip{u}{v}\defeq\int_a^b u^*(x)\,v(x)\,\rho(x)\,\dd x,\qquad
\norm{u}^2\defeq\ip{u}{u}=\int_a^b\lvert u\rvert^2\rho\,\dd x,\label{eq:prod}\\
L^2_\rho(a,b)\defeq\{u:\norm{u}<\infty\}/\!\sim,\qquad
u\sim v\iff\norm{u-v}=0\iff u=v\ \text{c.t.p.}\notag
\end{gather}
Propiedades: $\ip{u}{v}=\ip{v}{u}^*$; $\ip{u}{\alpha v+\beta w}=\alpha\ip{u}{v}+\beta\ip{u}{w}$;
$\ip{u}{u}\ge0$, $=0\iff u\sim0$. $L^2_\rho$ es completo (espacio de Hilbert).
\end{definicion}

\subsection{Simetría y propiedades espectrales}

\begin{proposicion}[Identidad de Lagrange]
Para $u,v\in\mathcal C^2(a,b)$:
\begin{equation}\label{eq:lagrange}
\ip{v}{\Lop u}-\ip{\Lop v}{u}=\Bigl[\,p\,\bigl(u\,v^{*\prime}-v^{*}u'\bigr)\Bigr]_a^b .
\end{equation}
\end{proposicion}

\begin{demostracion}
\begin{enumerate}
\item $\rho\bigl(v^*\Lop u-u\,(\Lop v)^*\bigr)=-v^*(pu')'+u(pv^{*\prime})'$ \quad($q,\rho,p$ reales: los términos en $q$ se cancelan).
\item $\dv{}{x}\bigl[p(u v^{*\prime}-v^*u')\bigr]=u'pv^{*\prime}+u(pv^{*\prime})'-v^{*\prime}pu'-v^*(pu')'=u(pv^{*\prime})'-v^*(pu')'$.
\item Integrando en $[a,b]$ con \eqref{eq:prod} se obtiene \eqref{eq:lagrange}.
\end{enumerate}
\end{demostracion}

\begin{definicion}[Operador simétrico]
$\Lop$ es simétrico en $D_{\Lop}$ si $\ip{v}{\Lop u}=\ip{\Lop v}{u}$ $\forall u,v\in D_{\Lop}$
$\iff$ $\bigl[p\,(u v^{*\prime}-v^*u')\bigr]_a^b=0$.
\end{definicion}

Condiciones de contorno que anulan el término de borde:
\begin{itemize}
\item \emph{Separadas} (Dirichlet, Neumann, Robin): $\alpha_c y(c)+\beta_c y'(c)=0$, $c\in\{a,b\}$, $(\alpha_c,\beta_c)\in\R^2\setminus\{0\}$
\[
\Rightarrow\ \begin{pmatrix}u(c)&u'(c)\\ v^*(c)&v^{*\prime}(c)\end{pmatrix}\!\begin{pmatrix}\alpha_c\\ \beta_c\end{pmatrix}=0
\ \Rightarrow\ u v^{*\prime}-v^*u'\big|_{c}=0 .
\]
\item \emph{Periódicas}: $p(a)=p(b)$, $y(a)=y(b)$, $y'(a)=y'(b)$.
\item \emph{Singulares}: $p(a)=p(b)=0$, o bien $a,b=\mp\infty$ con $p\,(uv^{*\prime}-v^*u')\to0$.
\end{itemize}

\begin{teorema}[Propiedades espectrales]
Sea $\Lop$ simétrico. Entonces:
\begin{enumerate}[label=(\alph*)]
\item $\Lop y=\lambda y\ \Rightarrow\ \lambda\in\R$.
\item $\Lop y_n=\lambda_n y_n$, $\Lop y_m=\lambda_m y_m$, $\lambda_n\neq\lambda_m\ \Rightarrow\ \ip{y_m}{y_n}=0$.
\item (S--L regular) $\lambda_0<\lambda_1<\dots\to\infty$ y $\{y_n\}$ es base ortogonal completa de $L^2_\rho$:
\begin{equation}\label{eq:desarrollo}
f=\sum_{n}c_n\,y_n,\qquad c_n=\frac{\ip{y_n}{f}}{\norm{y_n}^2},\qquad
\norm{f}^2=\sum_n\lvert c_n\rvert^2\norm{y_n}^2\ \ (\text{Parseval}).
\end{equation}
\end{enumerate}
\end{teorema}

\begin{demostracion}
\begin{enumerate}
\item[(a)] $\lambda\ip{y}{y}=\ip{y}{\Lop y}=\ip{\Lop y}{y}=\lambda^*\ip{y}{y}$, $\ \ip{y}{y}>0\Rightarrow\lambda=\lambda^*$.
\item[(b)] \[
\ip{y_m}{\Lop y_n}=\lambda_n\ip{y_m}{y_n},\qquad
\ip{\Lop y_m}{y_n}=\lambda_m^*\ip{y_m}{y_n}\overset{(a)}{=}\lambda_m\ip{y_m}{y_n}.
\]
Restando: $(\lambda_n-\lambda_m)\ip{y_m}{y_n}=0$.
\item[(c)] Completitud: se admite. Coeficientes: $\ip{y_m}{f}=\sum_n c_n\ip{y_m}{y_n}\overset{(b)}{=}c_m\norm{y_m}^2$.
\end{enumerate}
\end{demostracion}

\begin{observacion}[Analogía con $\R^3$]
Con $\varphi_n\defeq y_n/\norm{y_n}$:
\begin{center}
\renewcommand{\arraystretch}{1.35}
\begin{tabular}{>{\color{muted}}l c c}
\toprule
 & $\R^3$ & $L^2_\rho$\\
\midrule
base ortonormal & $\hat{\bm e}_i\cdot\hat{\bm e}_j=\delta_{ij}$ & $\ip{\varphi_n}{\varphi_m}=\delta_{nm}$\\
desarrollo & $\bm v=\sum_i(\hat{\bm e}_i\cdot\bm v)\,\hat{\bm e}_i$ & $f=\sum_n\ip{\varphi_n}{f}\varphi_n$\\
norma & $\lVert\bm v\rVert^2=\sum_i v_i^2$ & $\norm{f}^2=\sum_n\lvert\ip{\varphi_n}{f}\rvert^2$\\
\bottomrule
\end{tabular}
\end{center}
con $\delta_{ij}\defeq1$ si $i=j$, $0$ si $i\neq j$.
\end{observacion}

\begin{ejemplo}[Cuerda]
$\rho=p=1$, $q=0$, Dirichlet en $[0,\pi]$: $\Lop=-\dd^2/\dd x^2$ simétrico; $\lambda_n=n^2\in\R$;
$\ip{\sin nx}{\sin mx}=\tfrac{\pi}{2}\delta_{nm}$ $\Rightarrow$ \eqref{eq:fourier} es un caso de \eqref{eq:desarrollo}.
\end{ejemplo}

% =====================================================================
\section{Ecuación de Hermite}
% =====================================================================

\subsection{Forma de Sturm--Liouville}

\begin{definicion}[Ecuación de Hermite]
\begin{equation}\label{eq:hermite}
y''-2x\,y'+2\alpha\,y=0,\qquad x\in\R,\qquad \lambda\defeq2\alpha .
\end{equation}
\end{definicion}

\begin{proposicion}[Forma autoadjunta]
\begin{equation}\label{eq:autoadj}
y''+P(x)\,y'=\frac{1}{\rho}\,(\rho\,y')',\qquad \rho(x)=\exp\!\int^x P(s)\,\dd s .
\end{equation}
\end{proposicion}

\begin{demostracion}
$\frac{1}{\rho}(\rho y')'=y''+\frac{\rho'}{\rho}\,y'$, $\;\frac{\rho'}{\rho}=\dv{}{x}\ln\rho=P$.
\end{demostracion}

\begin{corolario}[Hermite como problema de S--L]
Con $P=-2x$: $\rho=\e^{-x^2}$, y \eqref{eq:hermite} $\iff\Lop y=\lambda y$ con
\begin{equation}\label{eq:hermiteSL}
\Lop y=-\e^{x^2}\bigl(\e^{-x^2}y'\bigr)'\!,\qquad
\rho=p=\e^{-x^2},\quad q=0,\quad \lambda=2\alpha,\quad D=\R .
\end{equation}
\end{corolario}

\begin{demostracion}
$-\e^{x^2}\bigl(\e^{-x^2}y'\bigr)'=-\e^{x^2}\e^{-x^2}\bigl(y''-2xy'\bigr)=-y''+2xy'\overset{\eqref{eq:hermite}}{=}2\alpha y$.
\end{demostracion}

\begin{observacion}[Simetría en $\R$]
Para $u,v$ de crecimiento polinómico (caso singular),
\[
\bigl[\e^{-x^2}(uv^{*\prime}-v^*u')\bigr]_{-\infty}^{\infty}=0\ \Rightarrow\ \Lop\ \text{simétrico}
\]
en $L^2(\R,\e^{-x^2}\dd x)$: $\lambda\in\R$ y autofunciones ortogonales.
En el manuscrito $p=-\e^{-x^2}$: equivale a $\tilde\Lop=-\Lop$, $\tilde\lambda=-2\alpha$; aquí $p>0$.
\end{observacion}

\subsection{Resolución mediante series}

$P(x)=-2x$, $Q(x)=2\alpha$ analíticas en $x=0$ $\Rightarrow$ $x=0$ punto regular (ordinario):
\begin{equation}
y=\sum_{n=0}^{\infty}a_nx^n,\qquad y'=\sum_{n=0}^{\infty}na_nx^{n-1},\qquad y''=\sum_{n=0}^{\infty}n(n-1)a_nx^{n-2}.
\end{equation}
Sustituyendo en \eqref{eq:hermite}:
\begin{equation}
\sum_{n=0}^{\infty}\Bigl[n(n-1)a_nx^{n-2}-2na_nx^n+2\alpha a_nx^n\Bigr]=0 .
\end{equation}
Reindexando ($n\to n+2$; los términos $n=0,1$ son nulos):
$\sum_{n=0}^\infty n(n-1)a_nx^{n-2}=\sum_{n=0}^\infty(n+2)(n+1)a_{n+2}x^n$, luego
\begin{equation}
\sum_{n=0}^{\infty}\Bigl[(n+2)(n+1)a_{n+2}+2(\alpha-n)a_n\Bigr]x^n=0\quad\forall x
\end{equation}
\begin{equation}\label{eq:recurrencia}
\Longrightarrow\quad(n+2)(n+1)\,a_{n+2}+2(\alpha-n)\,a_n=0
\quad\Longrightarrow\quad
\res{a_{n+2}=\frac{2(n-\alpha)}{(n+2)(n+1)}\,a_n .}
\end{equation}

Dos series independientes ($a_0$ genera la par, $a_1$ la impar): $y=a_0\,y_{\mathrm{par}}+a_1\,y_{\mathrm{impar}}$,
\begin{align}
y_{\mathrm{par}}&=1-\alpha x^2+\frac{\alpha(\alpha-2)}{6}x^4+\dots\notag\\
&=\sum_{k=0}^{\infty}\frac{(-2)^k\,\alpha(\alpha-2)\cdots(\alpha-2k+2)}{(2k)!}\,x^{2k},\label{eq:ypar}\\
y_{\mathrm{impar}}&=x+\frac{1-\alpha}{3}x^3+\frac{(3-\alpha)(1-\alpha)}{30}x^5+\dots\notag\\
&=\sum_{k=0}^{\infty}\frac{(-2)^k(\alpha-1)(\alpha-3)\cdots(\alpha-2k+1)}{(2k+1)!}\,x^{2k+1}.\label{eq:yimpar}
\end{align}

\begin{demostracion}[(coeficientes)]
De \eqref{eq:recurrencia}: $a_2=\frac{2(0-\alpha)}{2}a_0=-\alpha a_0$;
$a_4=\frac{2(2-\alpha)}{12}a_2=\frac{\alpha(\alpha-2)}{6}a_0$;
$a_3=\frac{2(1-\alpha)}{6}a_1=\frac{1-\alpha}{3}a_1$;
$a_5=\frac{2(3-\alpha)}{20}a_3=\frac{(3-\alpha)(1-\alpha)}{30}a_1$.
En general, $a_{2k}=a_0\prod_{j=0}^{k-1}\frac{2(2j-\alpha)}{(2j+2)(2j+1)}=a_0\,\frac{2^k\prod_{j=0}^{k-1}(2j-\alpha)}{(2k)!}$, y análogo para $a_{2k+1}$.
\end{demostracion}

\subsection{Convergencia y cuantización}

\begin{proposicion}[Comportamiento asintótico]
\begin{enumerate}[label=(\roman*)]
\item Radio de convergencia $R=\infty$.
\item Si la serie no se trunca: $y_{\mathrm{par}}\sim\e^{x^2}$, $\;y_{\mathrm{impar}}\sim x\,\e^{x^2}$ ($\lvert x\rvert\to\infty$) $\;\Rightarrow\; y\notin L^2_\rho(\R)$.
\end{enumerate}
\end{proposicion}

\begin{demostracion}
\begin{enumerate}
\item[(i)] $\displaystyle\Bigl\lvert\frac{a_{n+2}x^{n+2}}{a_nx^n}\Bigr\rvert=\frac{2\lvert n-\alpha\rvert}{(n+2)(n+1)}\,x^2\xrightarrow[n\to\infty]{}0\ \ \forall x$.
\item[(ii)] $n\gg\lvert\alpha\rvert$: $\dfrac{a_{n+2}}{a_n}\simeq\dfrac{2}{n}$. Para $\e^{x^2}=\sum_k\dfrac{x^{2k}}{k!}$, con $n=2k$:
$\dfrac{b_{n+2}}{b_n}=\dfrac{1}{k+1}=\dfrac{2}{n+2}\simeq\dfrac{2}{n}$. Mismo comportamiento de la cola $\Rightarrow$
\[
\norm{y}^2=\int_{-\infty}^{\infty}y^2\,\e^{-x^2}\dd x\sim\int^{\infty}\e^{2x^2}\e^{-x^2}\dd x=\int^{\infty}\e^{x^2}\dd x=\infty .
\]
\end{enumerate}
\end{demostracion}

\begin{teorema}[Cuantización]
Existe solución $y\in L^2_\rho(\R)$ de \eqref{eq:hermite} $\iff$ $\alpha=n\in\N$. Entonces $y$ es un polinomio de grado $n$ y paridad $(-1)^n$:
\begin{equation}
\res{\lambda_n=2n,\qquad y_n\propto H_n(x),\qquad n\in\N .}
\end{equation}
\end{teorema}

\begin{demostracion}
\begin{enumerate}
\item Por la proposición anterior, la serie debe truncarse: $\exists n$: $a_n\neq0$, $a_{n+2}=0$ $\overset{\eqref{eq:recurrencia}}{\iff}$ $\alpha=n$.
\item $n$ par: se trunca $y_{\mathrm{par}}$; $y_{\mathrm{impar}}$ no ($\alpha-(2j+1)\neq0$) $\Rightarrow a_1=0$.
\item $n$ impar: se trunca $y_{\mathrm{impar}}$; $y_{\mathrm{par}}$ no $\Rightarrow a_0=0$.
\end{enumerate}
\end{demostracion}

\begin{definicion}[Polinomios de Hermite]
$H_n$ es la solución polinómica de \eqref{eq:hermite} con $\alpha=n$, normalizada con coeficiente director $2^n$:
\begin{equation}\label{eq:explicita}
H_n(x)=\sum_{k=0}^{\lfloor n/2\rfloor}\frac{(-1)^k\,n!}{k!\,(n-2k)!}\,(2x)^{n-2k}.
\end{equation}
\end{definicion}

\begin{demostracion}[de \eqref{eq:explicita}]
Recurrencia descendente desde $a_n=2^n$: $a_{j}=\frac{(j+2)(j+1)}{2(j-n)}a_{j+2}$. Con $j=n-2k-2$:
\[
\frac{a_{n-2k-2}}{a_{n-2k}}=-\frac{(n-2k)(n-2k-1)}{4(k+1)},
\]
que coincide con el cociente de términos consecutivos ($k\to k+1$) de \eqref{eq:explicita}; $k=0$: $a_n=2^n$.
\end{demostracion}

\begin{center}
\renewcommand{\arraystretch}{1.3}
\begin{tabular}{>{\color{acc}}c l l}
\toprule
$n$ & $H_n(x)$ & desde \eqref{eq:ypar}--\eqref{eq:yimpar}\\
\midrule
0 & $1$ & $y_{\mathrm{par}}\big|_{\alpha=0}=1$\\
1 & $2x$ & $2\,y_{\mathrm{impar}}\big|_{\alpha=1}=2x$\\
2 & $4x^2-2$ & $-2\,y_{\mathrm{par}}\big|_{\alpha=2}=-2(1-2x^2)$\\
3 & $8x^3-12x$ & $-12\,y_{\mathrm{impar}}\big|_{\alpha=3}=-12\bigl(x-\tfrac23x^3\bigr)$\\
4 & $16x^4-48x^2+12$ & $12\,y_{\mathrm{par}}\big|_{\alpha=4}=12\bigl(1-4x^2+\tfrac43x^4\bigr)$\\
5 & $32x^5-160x^3+120x$ & $120\,y_{\mathrm{impar}}\big|_{\alpha=5}$\\
\bottomrule
\end{tabular}
\end{center}

\subsection{Oscilador armónico cuántico}

\begin{proposicion}[Cuantización de la energía]
\begin{gather}
-\frac{\hbar^2}{2m}\psi''+\frac12m\omega^2x^2\psi=E\psi,\qquad \psi\in L^2(\R)\notag\\
\Longrightarrow\quad
\res{E_n=\hbar\omega\bigl(n+\tfrac12\bigr),\qquad
\psi_n(\xi)=N_n\,H_n(\xi)\,\e^{-\xi^2/2},\qquad N_n=\Bigl(\frac{m\omega}{\pi\hbar}\Bigr)^{\!1/4}\frac{1}{\sqrt{2^nn!}}}
\end{gather}
con $\xi=\sqrt{m\omega/\hbar}\;x$.
\end{proposicion}

\begin{demostracion}
\begin{enumerate}
\item $\xi=\sqrt{m\omega/\hbar}\,x$, $\ \varepsilon\defeq2E/\hbar\omega$: $\ \psi_{\xi\xi}+(\varepsilon-\xi^2)\psi=0$.
\item $\psi=h(\xi)\,\e^{-\xi^2/2}$: $\ \psi_\xi=(h'-\xi h)\e^{-\xi^2/2}$, $\ \psi_{\xi\xi}=(h''-2\xi h'-h+\xi^2h)\e^{-\xi^2/2}$.
\item Sustituyendo: $h''-2\xi h'+(\varepsilon-1)h=0$ $\Rightarrow$ Hermite con $2\alpha=\varepsilon-1$.
\item $\norm{\psi}^2_{L^2}=\int h^2\e^{-\xi^2}\dd\xi=\norm{h}^2_{L^2_\rho}<\infty$ $\Rightarrow$ $\alpha=n$ $\Rightarrow$ $\varepsilon=2n+1$.
\item $N_n$: de \eqref{eq:ortog} y $\dd x=\sqrt{\hbar/m\omega}\,\dd\xi$.
\end{enumerate}
\end{demostracion}

\begin{figure}[ht]
\centering
\oscilador{0.82\linewidth}{left}
\caption{Potencial $V/\hbar\omega=\xi^2/2$ (azul claro) y funciones $\psi_n(\xi)\propto H_n(\xi)\,\e^{-\xi^2/2}$, $n=0,\dots,5$,
desplazadas a su nivel $E_n/\hbar\omega=n+\frac12$. El segmento tenue de cada nivel es la región clásicamente permitida
$\lvert\xi\rvert\le\sqrt{2n+1}$; $\psi_n$ tiene $n$ nodos y paridad $(-1)^n$.}
\label{fig:oscilador}
\end{figure}

% =====================================================================
\section{Fórmula de Rodrigues y ortogonalidad}
% =====================================================================

\begin{teorema}[Fórmula de Rodrigues]
\begin{equation}\label{eq:rodrigues}
\res{H_n(x)=(-1)^n\,\e^{x^2}\dvn{}{x}{n}\e^{-x^2}=\frac{(-1)^n}{\rho}\,\dvn{\rho}{x}{n},\qquad\rho=\e^{-x^2}.}
\end{equation}
Los $H_n$ están generados por el peso $\rho$ (una gaussiana).
\end{teorema}

\begin{demostracion}
Sea $w\defeq\e^{-x^2}$, $\ w^{(n)}\defeq\dd^nw/\dd x^n$.
\begin{enumerate}
\item $w'=-2xw\ \Rightarrow\ w'+2xw=0$.
\item Derivando $n+1$ veces (Leibniz: $(xw)^{(n+1)}=xw^{(n+1)}+(n+1)w^{(n)}$):
\[
w^{(n+2)}+2x\,w^{(n+1)}+2(n+1)\,w^{(n)}=0 .
\]
\item $u\defeq w^{(n)}\defeq(-1)^n\e^{-x^2}H$:
$\ u'=(-1)^n\e^{-x^2}(H'-2xH)$, $\ u''=(-1)^n\e^{-x^2}\bigl(H''-4xH'+(4x^2-2)H\bigr)$.
\item Sustituyendo en $u''+2xu'+2(n+1)u=0$:
\[
H''-4xH'+(4x^2-2)H+2xH'-4x^2H+2(n+1)H=H''-2xH'+2nH=0 .
\]
\item Inducción: $w^{(n)}=\bigl[(-2x)^n+\mathcal O(x^{n-2})\bigr]w\Rightarrow H=2^nx^n+\dots$ (polinomio, grado $n$, director $2^n$) $\Rightarrow H=H_n$.
\end{enumerate}
\end{demostracion}

\begin{ejemplo}
$w'=-2x\,w$, $\ w''=(4x^2-2)\,w$, $\ w'''=\bigl[8x-(4x^2-2)2x\bigr]w=(-8x^3+12x)\,w$
$\ \overset{\eqref{eq:rodrigues}}{\Longrightarrow}\ H_1=2x$, $H_2=4x^2-2$, $H_3=8x^3-12x$.
\end{ejemplo}

\begin{teorema}[Ortogonalidad]
\begin{equation}\label{eq:ortog}
\res{\ip{H_n}{H_m}=\int_{-\infty}^{\infty}H_n(x)\,H_m(x)\,\e^{-x^2}\dd x=2^n\,n!\,\sqrt\pi\;\delta_{nm}.}
\end{equation}
\end{teorema}

\begin{demostracion}
Sin pérdida de generalidad, $m\le n$.
\begin{enumerate}
\item Rodrigues: $\ip{H_m}{H_n}=(-1)^n\int_{-\infty}^{\infty}H_m\,w^{(n)}\,\dd x$.
\item $n$ integraciones por partes; términos de borde nulos:
\[
\bigl[H_m^{(k)}w^{(n-k-1)}\bigr]_{-\infty}^{\infty}=\bigl[(-1)^{n-k-1}H_m^{(k)}H_{n-k-1}\,\e^{-x^2}\bigr]_{-\infty}^{\infty}=0
\]
\[
\ip{H_m}{H_n}=(-1)^n(-1)^n\int_{-\infty}^{\infty}H_m^{(n)}\,\e^{-x^2}\dd x .
\]
\item $m<n$: $\deg H_m<n\Rightarrow H_m^{(n)}\equiv0\Rightarrow\ip{H_m}{H_n}=0$.
\item $m=n$: $H_n^{(n)}=2^n\,n!$ y $\int_{-\infty}^\infty\e^{-x^2}\dd x=\sqrt\pi$, pues
\[
\Bigl(\int_{-\infty}^{\infty}\e^{-x^2}\dd x\Bigr)^2=\iint_{\R^2}\e^{-(x^2+y^2)}\dd x\,\dd y=\int_0^{2\pi}\!\!\int_0^\infty\e^{-r^2}r\,\dd r\,\dd\theta=\pi .
\]
\end{enumerate}
\end{demostracion}

\begin{corolario}[Base ortonormal de $L^2(\R,\e^{-x^2}\dd x)$]
\begin{gather}
\varphi_n\defeq\frac{H_n}{\sqrt{2^nn!\sqrt\pi}},\qquad \ip{\varphi_n}{\varphi_m}=\delta_{nm},\notag\\
f=\sum_{n=0}^\infty c_nH_n,\quad c_n=\frac{1}{2^nn!\sqrt\pi}\int_{-\infty}^{\infty}f\,H_n\,\e^{-x^2}\dd x .
\end{gather}
\end{corolario}

\begin{ejemplo}[Desarrollo de $x^2$]
$c_0=\frac{1}{\sqrt\pi}\int x^2\e^{-x^2}\dd x=\frac{1}{\sqrt\pi}\cdot\frac{\sqrt\pi}{2}=\frac12$,
$\ c_2=\frac{1}{8\sqrt\pi}\int x^2(4x^2-2)\e^{-x^2}\dd x=\frac{1}{8\sqrt\pi}\bigl(4\cdot\tfrac{3\sqrt\pi}{4}-2\cdot\tfrac{\sqrt\pi}{2}\bigr)=\frac14$,
$\ c_{n}=0$ si $n\neq0,2$ $\Rightarrow$ $x^2=\tfrac12H_0+\tfrac14H_2$ \ (comprobación: $\tfrac12+\tfrac14(4x^2-2)=x^2$).
\end{ejemplo}

\begin{proposicion}[Función generatriz]
\begin{equation}\label{eq:generatriz}
g(x,t)\defeq\e^{2xt-t^2}=\sum_{n=0}^\infty H_n(x)\,\frac{t^n}{n!}.
\end{equation}
\end{proposicion}

\begin{demostracion}
$\e^{2xt-t^2}=\e^{x^2}\e^{-(x-t)^2}$; $\ \pd_t f(x-t)=-\pd_x f(x-t)\Rightarrow
\pd_t^n\e^{-(x-t)^2}\big|_{t=0}=(-1)^n\dvn{}{x}{n}\e^{-x^2}$. Taylor en $t$ y \eqref{eq:rodrigues}.
\end{demostracion}

% =====================================================================
\section{Operadores escalera}
% =====================================================================

\begin{definicion}[Operadores $A_\pm$]
Sobre polinomios ($\subset L^2_\rho(\R)$):
\begin{equation}
\Am\defeq\dv{}{x},\qquad \Ap\defeq-\dv{}{x}+2x .
\end{equation}
\end{definicion}

\begin{proposicion}[Factorización]
$\Ap\Am=-\dvn{}{x}{2}+2x\dv{}{x}=\Lop$, \ luego \eqref{eq:hermite} $\iff$
\begin{equation}
\Bigl(-\dv{}{x}+2x\Bigr)\dv{}{x}\,y_\alpha=2\alpha\,y_\alpha\iff\res{\Ap\Am\,y_\alpha=2\alpha\,y_\alpha .}
\end{equation}
\end{proposicion}

\begin{proposicion}[Adjunción]
$\ip{f}{\Am g}=\ip{\Ap f}{g}$ en $L^2_\rho(\R)$, es decir, $\Ap=\Am^{\dagger}$.
\end{proposicion}

\begin{demostracion}
$\int f^*g'\e^{-x^2}\dd x=\bigl[f^*g\,\e^{-x^2}\bigr]_{-\infty}^{\infty}-\int\bigl(f^*\e^{-x^2}\bigr)'g\,\dd x
=\int\bigl(-f^{*\prime}+2xf^*\bigr)g\,\e^{-x^2}\dd x=\ip{\Ap f}{g}$.
\end{demostracion}

\begin{corolario}[Cota inferior del espectro]
$2\alpha\norm{y_\alpha}^2=\ip{y_\alpha}{\Ap\Am y_\alpha}=\ip{\Am y_\alpha}{\Am y_\alpha}=\norm{y_\alpha'}^2\ge0$
$\Rightarrow\alpha\ge0$; $\ \alpha=0\iff y'=0\iff y=\mathrm{const}\propto H_0$.
\end{corolario}

\begin{proposicion}[Conmutador]
\begin{equation}
[\Am,\Ap]\defeq\Am\Ap-\Ap\Am=2\qquad\bigl(\Am\Ap=\Ap\Am+2\bigr).
\end{equation}
Si conmutan, $[\Am,\Ap]=0$ (\emph{conmutador}, como para dos matrices).
\end{proposicion}

\begin{demostracion}
\[
[\Am,\Ap]\,y=\dv{}{x}\Bigl(-\dv{y}{x}+2xy\Bigr)-\Bigl(-\dv{}{x}+2x\Bigr)\dv{y}{x}
=-y''+2y+2xy'+y''-2xy'=2y\quad\forall y .
\]
A nivel de generadores (identidad de operadores): $[\Am,\Ap]=2$.
\end{demostracion}

\begin{teorema}[Operadores escalera]
Si $\Ap\Am\,y_\alpha=2\alpha\,y_\alpha$:
\begin{equation}
\Ap\Am\,(\Ap y_\alpha)=2(\alpha+1)\,\Ap y_\alpha,\qquad
\Ap\Am\,(\Am y_\alpha)=2(\alpha-1)\,\Am y_\alpha .
\end{equation}
$\Ap$ sube ($\alpha\to\alpha+1$) y $\Am$ baja ($\alpha\to\alpha-1$).
\end{teorema}

\begin{demostracion}
\begin{align*}
\Ap\Am(\Ap y_\alpha)&=\Ap(\Am\Ap y_\alpha)=\Ap\bigl[(\Ap\Am+2)y_\alpha\bigr]=\Ap\bigl[2\alpha y_\alpha+2y_\alpha\bigr]=2(\alpha+1)\Ap y_\alpha,\\
\Ap\Am(\Am y_\alpha)&=(\Am\Ap-2)\Am y_\alpha=\Am(\Ap\Am y_\alpha)-2\Am y_\alpha=2(\alpha-1)\Am y_\alpha .
\end{align*}
\end{demostracion}

\begin{corolario}[Cuantización algebraica]
$\alpha\in\N$.
\end{corolario}

\begin{demostracion}
$\Am^jy_\alpha\neq0\Rightarrow$ autofunción con $\alpha-j\ge0$ (cota inferior). Luego $\exists k$: $\Am^{k}y_\alpha=0$, $\Am^{k-1}y_\alpha\neq0$;
$\ 0=\Ap\Am(\Am^{k-1}y_\alpha)=2(\alpha-k+1)\Am^{k-1}y_\alpha\Rightarrow\alpha=k-1\in\N$.
\end{demostracion}

\begin{proposicion}[Acción sobre $H_n$]
\begin{equation}\label{eq:escalera}
\res{\Ap H_n=H_{n+1},\qquad \Am H_n=H_n'=2n\,H_{n-1}.}
\end{equation}
\end{proposicion}

\begin{demostracion}
\begin{enumerate}
\item $w^{(n)}=(-1)^nH_nw\Rightarrow w^{(n+1)}=(-1)^n(H_n'-2xH_n)w$;
\eqref{eq:rodrigues}: $H_{n+1}=(-1)^{n+1}\frac{w^{(n+1)}}{w}=2xH_n-H_n'=\Ap H_n$.
\item $\Am H_n$: polinomio de grado $n-1$, autofunción con $\alpha=n-1$ $\Rightarrow\Am H_n=c_nH_{n-1}$.
\item $2nH_n=\Ap\Am H_n=c_n\Ap H_{n-1}=c_nH_n\Rightarrow c_n=2n$.
\end{enumerate}
\end{demostracion}

\begin{corolario}[Recurrencias]
\begin{equation}
H_n=\Ap^{\,n}H_0=\Bigl(2x-\dv{}{x}\Bigr)^{\!n}1,\qquad
H_{n+1}=2x\,H_n-2n\,H_{n-1},\qquad H_n'=2n\,H_{n-1}.
\end{equation}
\end{corolario}

\begin{ejemplo}
$H_2=\Ap(2x)=-2+4x^2$;
$\ H_3=2x(4x^2-2)-4\cdot2x=8x^3-12x$;
$\ H_4=2x(8x^3-12x)-6(4x^2-2)=16x^4-48x^2+12$;
$\ H_5=2xH_4-8H_3=32x^5-160x^3+120x$.
\end{ejemplo}

\begin{figure}[ht]
\centering
\begin{tikzpicture}[x=1cm,y=0.95cm,font=\small]
  \foreach \n/\col in {0/c0,1/c1,2/c2,3/c3,4/c4}{
    \draw[\col, line width=0.9pt] (0,\n) -- (4.2,\n);
    \node[anchor=east,\col] at (-0.15,\n) {$H_{\n}$};
    \node[anchor=west,muted] at (4.35,\n) {$\Ap\Am H_{\n}=\pgfmathparse{int(2*\n)}\pgfmathresult\,H_{\n}$};
  }
  \foreach \n in {0,1,2,3}{
    \draw[-{Stealth[length=2.2mm]}, acc, line width=0.7pt] (1.1,\n+0.08) -- (1.1,\n+0.92)
      node[midway,left,font=\scriptsize] {$\Ap$};
    \pgfmathtruncatemacro{\m}{\n+1}
    \pgfmathtruncatemacro{\k}{2*(\n+1)}
    \draw[-{Stealth[length=2.2mm]}, gold, line width=0.7pt] (3.1,\n+0.92) -- (3.1,\n+0.08)
      node[midway,right,font=\scriptsize] {$\Am:\ \times\k$};
  }
  \draw[-{Stealth[length=2.2mm]}, gold, line width=0.7pt] (3.1,-0.08) -- (3.1,-0.75);
  \node[gold,anchor=west,font=\scriptsize] at (3.2,-0.75) {$\Am H_0=0$};
  \node[muted,font=\scriptsize] at (2.1,4.55) {$\vdots$};
\end{tikzpicture}
\caption{Escalera de autofunciones de $\Lop=\Ap\Am$: $\Ap H_n=H_{n+1}$ (azul), $\Am H_n=2nH_{n-1}$ (dorado).
El espectro $\{2n\}_{n\in\N}$ está acotado inferiormente y $\Am$ aniquila el estado fundamental $H_0$.}
\label{fig:escalera}
\end{figure}

\begin{observacion}[Relación con $a,a^\dagger$]
Conjugando con $\e^{\xi^2/2}$ ($\psi=\e^{-\xi^2/2}h$):
\begin{gather}
\e^{-\xi^2/2}\,\Am\,\e^{\xi^2/2}=\dv{}{\xi}+\xi=\sqrt2\,a,\qquad
\e^{-\xi^2/2}\,\Ap\,\e^{\xi^2/2}=-\dv{}{\xi}+\xi=\sqrt2\,a^\dagger,\notag\\
[\Am,\Ap]=2\iff[a,a^\dagger]=1,\qquad
a^\dagger a=\tfrac12\Bigl(\xi^2-\dvn{}{\xi}{2}-1\Bigr),\notag\\
\hat H=\frac{\hbar\omega}{2}\Bigl(-\dvn{}{\xi}{2}+\xi^2\Bigr)\ \Rightarrow\ \hat H=\hbar\omega\bigl(a^\dagger a+\tfrac12\bigr),\quad a^\dagger a\,\psi_n=n\,\psi_n .
\end{gather}
\end{observacion}

\begin{metodo}{obtención de $H_n$}
\begin{enumerate}
\item \emph{Series}: \eqref{eq:recurrencia} con $\alpha=n$, truncar y fijar el director $2^n$.
\item \emph{Rodrigues}: \eqref{eq:rodrigues}, derivando $\e^{-x^2}$ $n$ veces.
\item \emph{Escalera}: $H_{n+1}=\Ap H_n=2xH_n-H_n'$, o bien $H_{n+1}=2xH_n-2nH_{n-1}$ desde $H_0=1$, $H_1=2x$.
\item \emph{Generatriz}: coeficiente de $t^n/n!$ en $\e^{2xt-t^2}$.
\item \emph{Comprobar}: paridad $(-1)^n$, $n$ raíces reales, $\ip{H_n}{H_n}=2^nn!\sqrt\pi$.
\end{enumerate}
\end{metodo}

% =====================================================================
\clearpage
\section{Formulario}
% =====================================================================
\newcommand{\fh}[1]{\par\smallskip{\sffamily\bfseries\color{acc}#1}\ \ }
\newcommand{\fhl}[1]{\par\smallskip{\sffamily\bfseries\color{lav}#1}\ \ }
\newcommand{\fhg}[1]{\par\smallskip{\sffamily\bfseries\color{gold}#1}\ \ }
{\footnotesize
\setlength{\columnsep}{1.4em}
\setlength{\columnseprule}{0.4pt}
\renewcommand{\columnseprulecolor}{\color{rulec}}
\setlength{\parskip}{0.35ex}
\begin{multicols}{2}
\fh{Separación} $y=X(x)T(t)$ en $k\,\pd_t^2y=\pd_x(u\pd_xy)-My$:
\[\ddot T=-\lambda T,\ \ T=c_\pm\e^{\pm\iu\sqrt\lambda t};\]\[\tfrac1k\bigl[-(uX')'+MX\bigr]=\lambda X.\]
Cuerda fija en $[0,\pi]$: $\lambda_n=n^2$, $X_n=\sin nx$, $\ip{X_n}{X_m}=\frac\pi2\delta_{nm}$.

\fhl{Sturm--Liouville} $\Lop y=\frac1\rho\bigl[-(py')'+qy\bigr]=\lambda y$, $\rho,p>0$.
\[\ip{u}{v}=\int_a^bu^*v\,\rho\,\dd x,\quad \norm{u}^2=\int_a^b\lvert u\rvert^2\rho\,\dd x .\]
Lagrange: $\ip{v}{\Lop u}-\ip{\Lop v}{u}=\bigl[p(uv^{*\prime}-v^*u')\bigr]_a^b$.\\
Simétrico $\Rightarrow\lambda\in\R$, $(\lambda_n-\lambda_m)\ip{y_m}{y_n}=0$.\\
$f=\sum c_ny_n$, $\ c_n=\ip{y_n}{f}/\norm{y_n}^2$; Parseval $\norm f^2=\sum\lvert c_n\rvert^2\norm{y_n}^2$.\\
Forma autoadjunta: $y''+Py'=\frac1\rho(\rho y')'$, $\rho=\e^{\int P}$.

\fhg{Hermite} $y''-2xy'+2\alpha y=0$; $\ \rho=p=\e^{-x^2}$, $q=0$, $\lambda=2\alpha$.
\[a_{n+2}=\frac{2(n-\alpha)}{(n+2)(n+1)}\,a_n,\quad y=a_0y_{\mathrm{par}}+a_1y_{\mathrm{impar}}.\]
No truncada: $y\sim\e^{x^2}\notin L^2_\rho$ $\Rightarrow$ $\alpha=n\in\N$, $\lambda_n=2n$.
\[H_n=\sum_{k=0}^{\lfloor n/2\rfloor}\frac{(-1)^kn!}{k!(n-2k)!}(2x)^{n-2k}.\]
$H_0=1$, $H_1=2x$, $H_2=4x^2-2$, $H_3=8x^3-12x$,\\
$H_4=16x^4-48x^2+12$, $H_5=32x^5-160x^3+120x$.\\
$H_n(-x)=(-1)^nH_n(x)$; $\ H_n^{(n)}=2^nn!$.

\fh{Rodrigues y generatriz}
\[H_n=(-1)^n\e^{x^2}\frac{\dd^n}{\dd x^n}\e^{-x^2},\quad \e^{2xt-t^2}=\sum_{n\ge0}H_n\frac{t^n}{n!}.\]
\fh{Ortogonalidad}
\[\int_{-\infty}^\infty H_nH_m\e^{-x^2}\dd x=2^nn!\sqrt\pi\,\delta_{nm}\]
$\varphi_n=H_n/\sqrt{2^nn!\sqrt\pi}$; $\ c_n=\frac{1}{2^nn!\sqrt\pi}\int fH_n\e^{-x^2}\dd x$.\\
$\int_\R\e^{-x^2}=\sqrt\pi$, $\ \int_\R x^2\e^{-x^2}=\frac{\sqrt\pi}2$, $\ \int_\R x^4\e^{-x^2}=\frac{3\sqrt\pi}4$.

\fhl{Escalera} $\Am=\frac{\dd}{\dd x}$, $\ \Ap=-\frac{\dd}{\dd x}+2x=\Am^\dagger$ (en $L^2_\rho$).
\[\Ap\Am y_\alpha=2\alpha y_\alpha,\quad[\Am,\Ap]=2.\]
$\Ap y_\alpha\propto y_{\alpha+1}$, $\ \Am y_\alpha\propto y_{\alpha-1}$; $\ 2\alpha\norm{y}^2=\norm{y'}^2\ge0$.
\[\Ap H_n=H_{n+1},\quad H_n'=2nH_{n-1},\]\[H_{n+1}=2xH_n-2nH_{n-1}.\]
$H_n=(2x-\frac{\dd}{\dd x})^n1$.

\fhg{Oscilador cuántico} $\xi=\sqrt{m\omega/\hbar}\,x$, $\ \psi=h\,\e^{-\xi^2/2}$:
\[h''-2\xi h'+(\varepsilon-1)h=0,\ \ \varepsilon=\tfrac{2E}{\hbar\omega}\]\[\Rightarrow\ E_n=\hbar\omega\bigl(n+\tfrac12\bigr).\]
$\psi_n=\bigl(\frac{m\omega}{\pi\hbar}\bigr)^{1/4}\frac{1}{\sqrt{2^nn!}}H_n(\xi)\e^{-\xi^2/2}$;\\
$\sqrt2a=\xi+\pd_\xi$, $\ \sqrt2a^\dagger=\xi-\pd_\xi$, $\ [a,a^\dagger]=1$, $\ \hat H=\hbar\omega(a^\dagger a+\frac12)$.
\end{multicols}
}

\end{document}
